% Direct-authority diplomatic French for printed p.216.
\clearpage
\SourcePage{216}

\LayerTarget{tohoku:proof:coinduced-hom-acyclicity}
\LayerTarget{tohoku:definition:coinduced-module-b-prime}
\noindent y ferons opérer $G$ et $O$ par

\LayerTarget{tohoku:equation:coinduced-u-actions}
\[
(g'f)(g)=f(gg'),
\qquad
(\lambda f)(g)=\lambda^g f(g),
\]
\LayerTarget{tohoku:construction:coinduced-u-module-b-prime}
 on constate aussitôt que $B'$ devient ainsi un $U$-module, et qu'on
définit un $U$-homomorphisme injectif $\varphi:B\to B'$ en posant

\LayerTarget{tohoku:equation:coinduced-u-embedding}
\[
\varphi(b)(g)=g\mathbin{\cdot}b
\]
\LayerTarget{tohoku:reduction:hom-direct-summand-acyclicity}
Ainsi $B$ se plonge dans le $U$-module $B'$, et comme $B$ est injectif, c'est
un facteur direct de $B'$, donc $\operatorname{Hom}_O(A,B)$ est (en tant que
$G$-module) un facteur direct de $\operatorname{Hom}_O(A,B')$, et il suffit de
montrer que ce dernier est $\Gamma^G$-acyclique.
\LayerTarget{tohoku:definition:conjugation-action-on-hom}
Posons $H=\operatorname{Hom}_O(A,B)$, considérons $H$ comme un $G$-module en y
faisant opérer $G$ par
$u^g(a)=g\,u(g^{-1}a)$;
on vérifie alors immédiatement qu'on a un isomorphisme canonique

\LayerTarget{tohoku:equation:hom-coinduction-isomorphism}
\[
\operatorname{Hom}_O(A,B')
=\operatorname{Hom}_O
  \bigl(A,\operatorname{Hom}_{\mathrm Z}(\mathbf Z(G),B)\bigr)
=\operatorname{Hom}_{\mathrm Z}
  \bigl(\mathbf Z(G),\operatorname{Hom}_O(A,B)\bigr)
\]
\LayerTarget{tohoku:construction:two-coinduced-g-actions}
 en convenant de mettre sur le dernier membre
$H''=\operatorname{Hom}_G(\mathbf Z(G),H)$ la structure de $G$-module définie
par
$f^{g'}(g)=g'f(gg')$.
On écrit pour ceci, pour toute
$u\in\operatorname{Hom}_O
\bigl(A,\operatorname{Hom}_{\mathrm Z}(\mathbf Z(G),B)\bigr)$,
$u(a)(g)=g\mathbin{\cdot}u_g(a)$, alors les $u_g$ sont des
$O$-homomorphismes de $A$ dans $B$ et $u$ s'identifie à la famille des $u_g$,
identifiée à un élément de $\operatorname{Hom}_{\mathrm Z}(\mathbf Z(G),H)$.
Considérons aussi le $G$-module $H'$ obtenu en mettant sur
$\operatorname{Hom}_{\mathrm Z}(\mathbf Z(G),H)$ la structure de $G$-module
définie par
$(g'f)(g)=f(gg')$.
\LayerTarget{tohoku:equation:coinduced-action-intertwiner}
On obtient un $G$-isomorphisme $\psi$ de $H'$ sur $H''$ eu posant
$\psi(f)(g)=g f(g)$.
\LayerTarget{tohoku:conclusion:coinduced-hom-acyclicity}
Or il est bien connu que $H'$ est $\Gamma^G$-acyclique, quel que soit le
$G$-module $H$, il en est donc de même de $H''$, cqfd.

\LayerTarget{tohoku:spectral-sequence:three-equivariant-ext-sequences}
Le lemme 5.6.1 nous permet de tirer des formules (5.6.5) trois suites
spectrales aboutissant à $\operatorname{Ext}_{{\mathbf O},G}(X;A,B)$. Pour expliciter
les termes initiaux des deux dernières, il reste à expliciter les foncteurs
dérivés de $h_{{\mathbf O},A}$ et $\mathbf h_{{\mathbf O},A}$ considérés comme foncteurs sur
${\mathbf C}^{{\mathbf O}(G)}$.
\LayerTarget{tohoku:relation:derived-hom-functors}
Quand on les considère comme foncteurs sur la catégorie ${\mathbf C}^{\mathbf O}$ des ${\mathbf O}$-Modules
(sans opérateurs) leurs foncteurs dérivés sont par définition (4.2) les
foncteurs $\operatorname{Ext}_{\mathbf O}^n(X;A,B)$ resp. les foncteurs
$\mathbf{Ext}_{\mathbf O}^n(A\ B)$. Il en est encore ainsi quand on les considère comme
des foncteurs sur ${\mathbf C}^{{\mathbf O}(G)}$, comme il résulte aussitôt des définitions et du

\LayerTarget{tohoku:lemma:5.6.2}
\textsc{Lemme 5.6.2.}\quad \textit{Si $B$ est un $G$-${\mathbf O}$-Module
injectif, c'est un ${\mathbf O}$-Module injectif.}

\LayerTarget{tohoku:proof:lemma-5.6.2}
Comme dans le lemme 5.6.1, on est ramené au cas où $B$ est le faisceau produit
défini par une famille $(B_x)_{x\in X}$ de $U_x$-modules injectifs $B_x$. Il
suffit donc de prouver que les $B_x$ sont aussi des ${\mathbf O}(x)$-modules injectifs,
ce qui nous ramène au cas où $X$ est réduit à un point $x$.
\LayerTarget{tohoku:reduction:equivariant-injective-underlying}
En fait, la démonstration ne semble pas plus simple dans ce cas, et celle que
nous allons donner s'appliquerait en fait chaque fois qu'on a une catégorie
abélienne ${\mathbf C}$ (telle que ${\mathbf C}^{\mathbf O}$) où ``opère'' un groupe $G$, comme
dans\textsuperscript{9)}, permettant de passer à la catégorie ${\mathbf C}(G)$ (ici
${\mathbf C}^{{\mathbf O}(G)}$) des ``objets avec opérateur'' correspondante ; on suppose que ${\mathbf C}$
satisfait $AB$ 5) et admet un générateur : \textit{alors tout objet injectif $A$
de ${\mathbf C}(G)$ est aussi injectif dans ${\mathbf C}$}.
\LayerTarget{tohoku:construction:equivariant-injective-envelope}
Pour le prouver, il suffit évidemment de prouver que tout objet $A$ de ${\mathbf C}(G)$
se plonge dans un objet $M$ de ${\mathbf C}(G)$ qui est injectif dans ${\mathbf C}$. Pour ceci, il
suffit d'examiner la construction donnée dans 1.10 pour l'immersion de $A$
(considéré comme objet de ${\mathbf C}$) dans un objet injectif $M$ de ${\mathbf C}$ : on plonge
$A$

\TohokuContinuation{217}
